% SPDX-License-Identifier: CC-BY-4.0
\documentclass[11pt,letterpaper]{article}

\usepackage[T1]{fontenc}
\usepackage{amsmath,amssymb}
\usepackage{tikz}
\usepackage{pgfplots}
\pgfplotsset{compat=1.18}
\usepackage{algorithm}
\usepackage{algpseudocode}
\usepackage{booktabs}
\usepackage{enumitem}
\usepackage[hidelinks]{hyperref}

\newcommand{\N}{\mathbb{N}}
\newcommand{\Z}{\mathbb{Z}}
\newcommand{\Q}{\mathbb{Q}}
\newcommand{\R}{\mathbb{R}}
\newcommand{\C}{\mathbb{C}}
\newcommand{\Map}[3]{#1 \colon #2 \to #3}
\newcommand{\Mod}[1]{#1\text{-}\mathrm{Mod}}
\newcommand{\GroupHom}{\mathrm{Hom}_{\mathbf{Grp}}}
\newcommand{\RingHom}{\mathrm{Hom}_{\mathbf{Ring}}}

\newtheorem{theorem}{Theorem}
\newtheorem{lemma}[theorem]{Lemma}
\newtheorem{corollary}[theorem]{Corollary}
\newtheorem{definition}[theorem]{Definition}
\newtheorem{remark}[theorem]{Remark}
\newenvironment{proof}{\noindent\textit{Proof.}}{\hfill\(\square\)}

\title{Generalized Collatz Library}
\author{}
\date{\today}

\begin{document}
\maketitle
\section*{Entry into the residue lattice before division}

\begin{abstract}
If \(D\) does not divide \(c\), \(\gcd(c,D)=1\), and \(c\) divides \(q^k\), then every integer seed reaches the residue lattice \(c\Z\) in \(k\) steps of the affine map \(n\mapsto qn+c\). Those steps use no division. Entry into \(c\Z\) is not capture of the state \(1\), and it is not the assertion that the classical raw map reaches \(1\).
\end{abstract}

\section{Affine steps before division}

\begin{definition}
\label{def:affine}
For \(q,c\in\Z\), the affine map is \(\Map{A_{q,c}}{\Z}{\Z}\),
\[
A_{q,c}(n)=qn+c.
\]
Its iterates are \(A_{q,c}^{0}(n)=n\) and \(A_{q,c}^{j+1}(n)=A_{q,c}(A_{q,c}^{j}(n))\). Define weights by \(G_0(q)=0\) and \(G_{j+1}(q)=qG_j(q)+1\).
\end{definition}

The value \(A_{q,c}^{k}(n)\) is computed by \(k\) applications of \(n\mapsto qn+c\). Division by an integer \(D\) is not one of those applications.

\begin{lemma}
\label{lem:closed}
For all \(q,c,n\in\Z\) and every \(k\in\N\),
\[
A_{q,c}^{k}(n)=q^k n+cG_k(q).
\]
\end{lemma}

\begin{proof}
For \(k=0\), both sides equal \(n\), since \(q^0=1\) and \(G_0(q)=0\). If the identity holds at \(k\), then
\begin{align*}
A_{q,c}^{k+1}(n)
&=qA_{q,c}^{k}(n)+c\\
&=q(q^k n+cG_k(q))+c\\
&=q^{k+1}n+c(qG_k(q)+1)\\
&=q^{k+1}n+cG_{k+1}(q).
\end{align*}
Induction gives the identity at every \(k\).
\end{proof}

\begin{lemma}
\label{lem:power}
If \(c\mid q^k\), then \(c\mid A_{q,c}^{k}(n)\) for every \(n\in\Z\).
\end{lemma}

\begin{proof}
By Lemma~\ref{lem:closed}, \(A_{q,c}^{k}(n)=q^k n+cG_k(q)\). The hypothesis gives \(c\mid q^k n\). The second summand is \(c\) times an integer. A sum of multiples of \(c\) is a multiple of \(c\).
\end{proof}

\section{The offset lattice}

\begin{definition}
\label{def:hypotheses}
Integers \(q,c,D\) and a natural number \(k\) satisfy the lattice-entry hypotheses when
\[
D\nmid c,\qquad \gcd(c,D)=1,\qquad c\mid q^k.
\]
Here \(\gcd(c,D)\) denotes the least positive element of the ideal \(c\Z+D\Z\), or \(0\) if that ideal is \(\{0\}\).
\end{definition}

\begin{lemma}
\label{lem:transverse}
If \(\gcd(c,D)=1\), then for every \(m\in\Z\),
\[
D\mid cm\quad\Longleftrightarrow\quad D\mid m.
\]
\end{lemma}

\begin{proof}
If \(D\mid m\), then \(D\mid cm\). Conversely, the definition of the gcd and the hypothesis \(\gcd(c,D)=1\)
give integers \(x,y\) such that
\[
cx+Dy=1.
\]
Multiply by \(m\):
\[
m=(cx)m+(Dy)m.
\]
The second summand is divisible by \(D\). The hypothesis \(D\mid cm\) makes the first summand divisible by \(D\). Therefore \(D\mid m\).
\end{proof}

\begin{lemma}
\label{lem:offset-affine}
Let \(|D|\ge 2\) and suppose \(D\nmid c\). On the raw map that sends a divisible input to its quotient by \(D\) and every other input to \(qn+c\), the input \(c\) takes the affine branch:
\[
qc+c.
\]
\end{lemma}

\begin{proof}
The branch is affine precisely when \(D\) does not divide the input. That is the hypothesis at the input \(c\), and the affine value is \(qc+c\).
\end{proof}

\begin{theorem}
\label{thm:entry}
Assume the lattice-entry hypotheses. For every \(n\in\Z\),
\[
c\mid A_{q,c}^{k}(n).
\]
Thus every integer seed reaches the residue lattice \(c\Z\) at the \(k\)th step of \(n\mapsto qn+c\), hence within \(k\) affine steps, and that computation uses no division.
\end{theorem}

\begin{proof}
Lemma~\ref{lem:power} applies because \(c\mid q^k\), so \(c\) divides \(A_{q,c}^{k}(n)\). The state is the \(k\)-fold iterate of \(A_{q,c}\). Division by \(D\) is not used. The hypotheses \(D\nmid c\) and \(\gcd(c,D)=1\) are not used in this expansion. The first is the hypothesis of Lemma~\ref{lem:offset-affine}, and the second is the hypothesis of Lemma~\ref{lem:transverse}.
\end{proof}

\section{Entry is not capture}

For integers \(q,c,D\) with \(|D|\ge 2\), the raw map is
\[
F_{q,c,D}(n)=
\begin{cases}
n/D,&D\mid n,\\
qn+c,&D\nmid n.
\end{cases}
\]
Capture of a seed at the state \(1\) would mean that some raw iterate equals \(1\).

\begin{theorem}
\label{thm:not-capture}
There exist integers \(q,c,D\) with \(|D|\ge 2\), a natural number \(k\), and a seed \(n\in\Z\), satisfying the lattice-entry hypotheses, such that \(c\mid A_{q,c}^{k}(n)\) and, for every \(m\in\N\),
\[
F_{q,c,D}^{m}(n)\ne 1.
\]
Thus lattice entry is not capture of the state \(1\).
\end{theorem}

\begin{proof}
Take \(q=2\), \(c=2\), \(D=3\), \(k=1\), and \(n=2\). Then \(|D|=3\ge 2\). Three does not divide two, \(\gcd(2,3)=1\), and \(2\mid 2^1\). Lemma~\ref{lem:power} gives \(2\mid A_{2,2}(2)\). Explicitly, \(A_{2,2}(2)=6\).

Since \(3\nmid 2\), the raw step at \(2\) is affine, so \(F(2)=2\cdot 2+2=6\). Since \(3\mid 6\), the next raw step divides, so \(F(6)=6/3=2\). Induction on \(j\) gives \(F^{2j}(2)=2\). The case \(j=0\) is the seed. If the identity holds at \(j\), then
\[
F^{2(j+1)}(2)=F\bigl(F\bigl(F^{2j}(2)\bigr)\bigr)=F(F(2))=F(6)=2.
\]
One further raw step gives \(F^{2j+1}(2)=F\bigl(F^{2j}(2)\bigr)=F(2)=6\). Every natural number \(m\) has the form \(2j\) or \(2j+1\): the division identity \(m=(m/2)\cdot 2+(m\bmod 2)\) has remainder \(0\) or \(1\). Thus every raw iterate equals \(2\) or \(6\), and neither equals \(1\).
\end{proof}

\begin{remark}
\label{rem:classical}
For the classical parameters \((q,c,D)=(3,1,2)\), all lattice-entry
hypotheses hold for every \(k\), but the conclusion is only divisibility
by \(1\). This places no restriction on a state and supplies no hitting time. In particular, it does not assert that the classical raw map reaches \(1\). Theorem~\ref{thm:not-capture} separates lattice membership from that hitting statement for a different parameter triple. No orbit on \(c\Z\) is classified here, and the identity \(F_{q,c,D}(cn)=cF_{q,1,D}(n)\) is not claimed.
\end{remark}

\end{document}
